\section{Infra 与数据：规模化学习的底座}

上一章解释“模型为什么会变”，本章解释“模型为什么能放大”。当模型参数、数据和上下文都增长时，问题不再只是算法，而是显存、通信、并行、数据质量和集群稳定性。ZeRO 是 Zero Redundancy Optimizer，它通过分片优化器状态、梯度和参数来降低训练冗余；Scaling Law、Chinchilla、LAION、RefinedWeb 和 MegaScale 则构成了这一部分的其他主线。

\begin{figure}[H]
\centering
\includegraphics[width=0.86\textwidth]{infra-data-stack.png}
\caption{Infra 与数据栈：ZeRO、Scaling Law、开源数据和万卡集群共同放大模型。自制概念图，依据 01:52:58--02:21:29 对谈内容整理。}
\end{figure}

\begin{knowledgebox}{读图：底座不是幕后杂务}
没有 ZeRO 这类并行和状态分片，大模型放不进多 GPU；这里的状态包括 optimizer state，也就是 Adam/AdamW 的 \(m\)、\(v\) 等优化器状态。没有 Scaling Law 和 Chinchilla，训练预算不知道如何分配；没有 LAION、RefinedWeb 这类数据工程，模型没有足够可用语料；没有 MegaScale 这类集群工程，万卡训练就只是口号。
\end{knowledgebox}

\subsection{ZeRO：把状态切开}

本节先处理一个必须消化的系统概念。ZeRO 的核心是把训练中原本每张 GPU 都复制的状态切开，包括优化器状态、梯度和参数，从而降低单卡显存压力。它不是让模型变聪明的算法，但它让更大模型训练成为可能。

\begin{importantbox}{什么是 ZeRO：Zero Redundancy Optimizer}
ZeRO 是 Zero Redundancy Optimizer 的缩写，核心思想是把数据并行中的冗余状态做分片（sharding），让每张 GPU 只保存一部分训练状态，而不是完整复制。ZeRO-1 分片优化器状态，ZeRO-2 进一步分片梯度，ZeRO-3 进一步分片参数。这里的 optimizer state 通常指 Adam/AdamW 的一阶动量 \(m\) 和二阶动量 \(v\)，这些状态会显著增加训练显存。
\end{importantbox}

\subsection{Scaling Law 与 Chinchilla：预算分配的指挥棒}

Scaling Law 尝试描述模型性能如何随参数量、数据量和计算量变化。Chinchilla 则进一步提醒，在固定计算预算下，模型参数和训练 token 需要更合理配比；只堆参数而数据不足，可能不是最优。它们的意义不是给出永恒定律，而是让训练从拍脑袋变成可预测的预算分配问题。

\begin{knowledgebox}{Scaling Law 的直觉公式}
可以把损失的下降粗略理解为：
\[
L(N, D, C) \downarrow \quad \text{as} \quad N,D,C \uparrow
\]
其中，\(N\) 表示模型参数量，\(D\) 表示训练数据量，\(C\) 表示计算预算。这个式子不是精确公式，而是提醒我们：参数、数据和计算要配合增长，不能只押一个维度。
\end{knowledgebox}

\subsection{LAION、RefinedWeb 与开源数据}

上一节讨论预算如何分配，本节转向预算背后的燃料：数据。LAION-5B 体现了开源社区希望打破少数工业巨头垄断的努力；RefinedWeb 则说明互联网数据经过清洗、过滤和配比，仍然可以支撑强模型。数据不是越多越好，而是来源、质量、过滤、去重、版权和可复现性共同决定模型上限。

\begin{warningbox}{开源数据不是免费午餐}
开源数据降低了进入门槛，也带来质量、版权、偏见和安全问题。真正的数据工程不是下载一个大文件，而是持续做过滤、去重、标注、评测和追踪。
\end{warningbox}

\subsection{MegaScale：万卡训练的工程课}

数据决定模型能学什么，但要把数据真正喂给模型，还需要集群工程。本节的 MegaScale 代表大规模训练的工程现实。训练集群里 GPU 利用率、通信开销、故障恢复、调度、网络拓扑和内存带宽都会决定真实成本。讲者提到 MFU，即 Model FLOPs Utilization，反映模型实际使用硬件算力的比例；哪怕接近 50\% 已经很优秀，也意味着仍有大量算力被通信和等待吞掉。

\begin{knowledgebox}{术语消化：Infra 与数据}
\begin{tabularx}{\textwidth}{p{0.22\textwidth}p{0.34\textwidth}X}
\toprule
术语 & 解决的问题 & 为什么重要 \\
\midrule
ZeRO & 训练状态分片 & 让更大模型能放进多 GPU 显存。 \\
Scaling Law & 预测规模与性能关系 & 指导参数、数据和计算预算。 \\
Chinchilla & 数据/参数配比 & 防止只堆参数而训练 token 不足。 \\
LAION/RefinedWeb & 开源多模态/网页数据 & 让模型训练不完全依赖封闭巨头。 \\
MFU & 衡量硬件利用率 & 低 MFU 会让同样算力买来更少训练进度。 \\
\bottomrule
\end{tabularx}
\end{knowledgebox}

\subsection{本章小结}

Infra 与数据线说明，大模型不是只靠论文公式长大的。它需要显存分片、并行通信、规模规律、数据工程、开源社区和集群可靠性。模型范式定义方向，Infra 和数据决定方向能不能被规模化兑现。

